\documentclass[11pt,a4paper]{article}
\usepackage[T1]{fontenc}
\usepackage[utf8]{inputenc}
\usepackage[polish]{babel}
\usepackage{lmodern}
\usepackage{microtype}
\usepackage[a4paper,margin=2.3cm]{geometry}
\usepackage{amsmath,amssymb}
\usepackage{booktabs,longtable}
\usepackage{xcolor}
\usepackage{hyperref}
\usepackage{listings}
\hypersetup{colorlinks=true,linkcolor=blue,urlcolor=blue}
\title{Import historycznych danych PM2.5 GIOŚ\\
Diagnoza, niezmienniki i przepływ Bridge}
\author{Edward Szczypka\\\href{mailto:edward@szczypka.guru}{edward@szczypka.guru}}
\date{2026}
\begin{document}
\maketitle

\section{Wynik audytu}

Oficjalne roczne archiwa zostały pobrane poprawnie i przechodzą kontrolę CRC.
Parser odczytał następujące liczby poprawnych, nieujemnych wartości:

\begin{center}
\begin{tabular}{rrrr}
\toprule
Rok & Liczba serii & Poprawne wartości & Wiersze w SQLite\\
\midrule
2022 & 89 & 755956 & 0\\
2023 & 92 & 779380 & 0\\
2024 & 96 & 817648 & 0\\
\bottomrule
\end{tabular}
\end{center}

Oznacza to, że problem nie dotyczy dostępności danych GIOŚ. Awaria leżała
pomiędzy wyborem arkusza a warstwą persystencji.

\section{Zidentyfikowane błędy}

\subsection{Normalizacja nazwy zanieczyszczenia}

Kod aplikacji używa oznaczenia
\[
  \mathrm{PM}_{2.5},
\]
natomiast oficjalne pliki roczne mają nazwy typu
\texttt{2022\_PM25\_1g.xlsx}. Selektor musi utożsamiać warianty
\texttt{PM25}, \texttt{PM2.5} oraz \texttt{PM2,5}.

\subsection{Wiersz kodu stacji}

W arkuszu występuje wiersz porządkowy \texttt{Nr} zawierający wartości
$1,2,\ldots,n$ oraz osobny wiersz \texttt{Kod stacji}. Heurystyka oparta tylko
na liczbie unikalnych wartości wybierała wiersz porządkowy. Poprawiony algorytm
najpierw wyszukuje jawną etykietę \texttt{Kod stacji} i odrzuca sekwencje
kolejnych liczb naturalnych.

\section{Niezmiennik persystencji}

Niech $V$ oznacza liczbę poprawnych wartości w arkuszu, $I$ liczbę nowych
insertów, a $D$ liczbę rekordów pominiętych jako istniejące duplikaty.
Importer akceptuje arkusz wyłącznie wtedy, gdy
\[
  V = I + D.
\]
Jeżeli $V>0$ oraz $I+D=0$, wykonanie kończy się błędem i nie zapisuje markera
ukończenia w \texttt{state.json}.

\section{Przepływ Bridge}

\begin{verbatim}
GIOŚ ZIP
  -> HistoricalDataCacheBridge
       local | object_store | hybrid
  -> normalizacja PM25/PM2.5
  -> parser XLSX i Kod stacji
  -> SQLite batch insert
  -> walidacja V = I + D
  -> state.json
\end{verbatim}

Backend ObjectStore może być implementacją DigitalOcean Spaces, AWS S3,
MinIO albo lokalnego systemu plików. Wybór backendu nie zmienia kontraktu
naukowego parsera i persystencji.

\section{Progress i diagnostyka}

Dla każdej serii emitowane są: rok, parametr, kod stacji, liczba wykonanych
serii, liczba wszystkich serii, procent, nowe rekordy oraz duplikaty. Log jest
kierowany na standardowe wyjście i do pliku JSONL. Monitor PowerShell pokazuje
procent, ETA, PID, CPU i RAM.

\end{document}
